\section{Mathematical details for the audited margin vector}\label{app:neural}

\paragraph{Exploratory real-data runs.} The repository also holds a
state-held-out pilot of the proposal model on the public Census file and a later
comparison of linear, neural and quadratic proposers on the same file
(\texttt{scripts/realdata\_neural\_pilot.py} and
\texttt{scripts/realdata\_neural\_comparison.py}). Their linear baselines were
not trained to convergence and their differences lie within bootstrap intervals,
so they are not evidence for or against the procedure, and nothing in this paper
rests on them.

\paragraph{Experiment and estimand.}
Let $\mathcal D$ contain all development observations, candidate generation,
neural fitting, tuning, scenario definitions and proposed releases. Conditional
on $\mathcal D$, the shortlist $\Theta_0$, its cardinality $C$, and all
sample sizes $B_{\theta j}\geq1$ are fixed. An audit draw $Z_{\theta ji}$
includes the target, auxiliary data when freshly sampled, and any randomness
in the declared attack. The measurable success indicator
$S_{\theta ji}=s_{\theta j}(Z_{\theta ji})$ is one exactly when the declared
harm criterion is met. For fixed $(\theta,j)$ assume the draws are identically
distributed and independent conditional on $\mathcal D$, with law
$P_{\theta j}$. Thus
\begin{equation}
 r_j(\theta)=\int s_{\theta j}(z)\,dP_{\theta j}(z),\qquad
 K_{\theta j}=\sum_{i=1}^{B_{\theta j}}S_{\theta ji}\mid\mathcal D
 \sim\operatorname{Binom}(B_{\theta j},r_j(\theta)).
 \label{eq:neural-experiment}
\end{equation}
The risk is conditional on a frozen release and on the audit design. If target
records or auxiliary files are reused, the resulting dependent observations
must be treated as one sampling unit or justified by another valid confidence
procedure. Independence across distinct $(\theta,j)$ pairs is unnecessary.
For a finite population sampled without replacement, the binomial assumption
is inappropriate unless the attack design supplies independent replacement
draws; a hypergeometric or design-based bound is needed instead.

\paragraph{Exact inversion.}
Write $b=B_{\theta j}$, $k=K_{\theta j}$,
$\gamma=\alpha/(Cn)$, and
$F_p(k)=\sum_{t=0}^{k}\binom{b}{t}p^t(1-p)^{b-t}$. For $k<b$,
$F_p(k)$ is continuous and strictly decreasing in $p\in(0,1)$.
The upper endpoint $u(k)$ is the unique solution of $F_{u(k)}(k)=\gamma$;
set $u(b)=1$. The binomial--beta identity gives
\begin{equation}
 F_p(k)=1-I_p(k+1,b-k),\qquad
 u(k)=F^{-1}_{\operatorname{Beta}(k+1,b-k)}(1-\gamma).
 \label{eq:neural-inversion}
\end{equation}
For any fixed true $p$, the event $\{u(K)<p\}$ is a lower-tail set of
integer counts: by monotonicity it equals $\{F_p(K)<\gamma\}$, apart from
the harmless equality boundary. If $k_*$ is the largest count in this set,
$\Pr_p\{u(K)<p\}=F_p(k_*)<\gamma$; if the set is empty the probability is
zero. Hence $\Pr_p\{p\leq u(K)\}\geq1-\gamma$ for every $p\in[0,1]$.
The possible strict inequality reflects the discreteness of the count.

\paragraph{Familywise selection.}
Define the event
\begin{equation}
 E=\bigcap_{\theta\in\Theta_0}\bigcap_{j\in\mathcal J}
 \{r_j(\theta)\leq\overline r_{\theta j}\}.
 \label{eq:neural-good-event}
\end{equation}
Conditional on every development history for which the assumptions hold,
$\Pr(E^c\mid\mathcal D)\leq\sum_{\theta,j}\gamma=\alpha$.
Any measurable decision rule $\widehat\theta$ using the complete vector of
final bounds, including the choice of the greatest utility among accepted
candidates, satisfies the deterministic implication
\begin{equation}
 E\ \Longrightarrow\
 \bigl[\widehat\theta\text{ accepted}\bigr]\ \Longrightarrow\
 \max_{j\in\mathcal J}\frac{r_j(\widehat\theta)}{\tau_j}\leq1.
 \label{eq:neural-implication}
\end{equation}
The probability of a false certification is therefore at most $\alpha$
conditionally and, by $\Pr(E^c)=\E[\Pr(E^c\mid\mathcal D)]$, unconditionally.
This is a frequentist guarantee for the declared finite family; it is not a
posterior probability that a particular accepted release is safe. An audit
outcome used to add candidates, alter scenarios, retrain the network or select
sample sizes destroys the conditioning argument unless a separate sequential
or adaptive design is proved valid.

\paragraph{Budget and an illustrative calculation.}
At $k=0$, equation \eqref{eq:neural-inversion} simplifies to
$u(0)=1-\gamma^{1/b}$. Consequently $u(0)\leq\tau$ if and only if
$(1-\tau)^b\leq\gamma$, equivalently
\begin{equation}
 b\geq\left\lceil\frac{\log\gamma}{\log(1-\tau)}\right\rceil.
 \label{eq:neural-budget-proof}
\end{equation}
For illustration only, $C=4$, $n=2$, $\alpha=0.05$ and $\tau=0.12$ give
$\gamma=0.00625$ and $b=40$. The zero-success endpoint is approximately
$0.11916$ at $b=40$, but approximately $0.12202$ at $b=39$.
These numbers do not establish the risk of any real release. The total number
of attack trials is $\sum_{\theta,j}B_{\theta j}$; when each of $Cn$ pairs
uses the same $b$, it is $Cnb$. For small $\tau$,
$-\log(1-\tau)=\tau+O(\tau^2)$, so the zero-success requirement grows as
$\log(Cn/\alpha)/\tau$ per pair. The neural output dimension changes this
multiplicity even when its input feature dimension $d$ stays fixed.

\paragraph{Neural objective and its limits.}
Let $\pi_j(x)=\Pr(Y_j=1\mid X=x)$ in the development distribution. For a
positive weight $w_j(x)$ depending only on $x$ and $j$, the conditional
expected binary cross entropy at a predicted probability $q$ is
\begin{equation}
 \ell_x(q)=-w_j(x)[\pi_j(x)\log q+(1-\pi_j(x))\log(1-q)],
 \quad \ell_x'(q)=w_j(x)\frac{q-\pi_j(x)}{q(1-q)},
 \label{eq:neural-proper}
\end{equation}
with $\ell_x''(q)>0$ at its unique minimum $q=\pi_j(x)$ whenever
$0<\pi_j(x)<1$. Outcome-dependent weights instead change that minimiser.
The implemented one-hidden-layer ReLU model is fitted by ordinary
backpropagation of this objective \citep[ch.~6]{goodfellow2016}; the gradients are
textbook ones, and fitting does not make the predictions exact. Write $W_1$ and
$W_2$ for its fitted input-to-hidden and hidden-to-output weight matrices. Then
ReLU's Lipschitz constant of one and sigmoid's scalar derivative at most
$1/4$ give the predictor sensitivity bound
$\|\widehat{\boldsymbol r}(x)-\widehat{\boldsymbol r}(x')\|_2
\leq\|W_2\|_2\|W_1\|_2\|x-x'\|_2/4$.
It controls changes in \emph{predictions}, not errors relative to actual
risks, and is no substitute for the audit.

\paragraph{The input dimension.}
Suppose, as an additional assumption, the true conditional risk is
$L$-Lipschitz in $\|\cdot\|_\infty$ on $[0,1]^d$. A grid of spacing
$h\leq\epsilon/L$ lets a value at a nearby grid point approximate the
true risk to within $\epsilon$, requiring at most
$(1+\lceil L/\epsilon\rceil)^d$ grid points. Estimating those values adds
sampling cost at each point; the grid count alone already grows exponentially
in $d$. If no regularity links unsampled inputs to observed inputs, two risk
functions can agree on every observed feature vector and disagree arbitrarily
at an unobserved vector, so there is no distribution-free uniform interpolation
guarantee. Structural assumptions may help prediction, but the finite-family
audit guarantee above does not rely on them.

\paragraph{Distribution shift.}
Using $\mathrm{TV}(P,Q)=\sup_A|P(A)-Q(A)|$ on the complete audit-draw space,
the layer-cake representation for any success function $0\leq s\leq1$ yields
$|\E_Qs-\E_Ps|\leq\mathrm{TV}(P,Q)$. If a justified deterministic bound
$\mathrm{TV}(P_{\theta j},Q_{\theta j})\leq\delta_{\theta j}$ holds for
all inspected pairs, a simultaneous deployed-risk bound is
\begin{equation}
 \overline r^{\rm deploy}_{\theta j}
 =\min\{1,\overline r_{\theta j}+\delta_{\theta j}\},\qquad
 T^{\rm deploy}(\theta)=\min_j\left(1-
 \frac{\overline r^{\rm deploy}_{\theta j}}{\tau_j}\right).
 \label{eq:neural-shift}
\end{equation}
The same familywise probability applies because the TV inequality holds
pointwise on the simultaneous coverage event. A TV estimate, itself uncertain,
needs its own valid upper bound and error allocation. A distance between neural
feature embeddings alone cannot establish a TV bound on the full experiment.
